\documentclass[11pt]{article}
\usepackage[margin=1in]{geometry}
\usepackage{booktabs}
\usepackage[hidelinks]{hyperref}
\usepackage{natbib}

% Every number below comes from generated/results.tex, which `hof paper`
% writes from a pipeline run. Never type a result by hand.
\input{generated/results.tex}

\title{An Objective Ranking of ATP Players for the International Tennis Hall of Fame}
\author{Example Team}
\date{Draft}

\begin{document}
\maketitle

\begin{abstract}
We score every Open Era ATP career from match results and weekly rankings,
fit a Hall of Fame threshold on players eligible up to \SplitYear, and test it
on the \TestPlayers{} contenders who became eligible after that. This draft runs
on synthetic data; the numbers are placeholders for the method, not findings.
\end{abstract}

\section{Introduction}
The International Tennis Hall of Fame inducts players by vote~\citep{ithf}.
We ask whether a transparent score built from on-court results can reproduce
those votes, and where it disagrees.

\section{Data}
Match results, player records and weekly top-10 rankings, pinned to one data
snapshot. See \texttt{data/snapshot.toml} in the repository.

\section{Method}
A player's career score is a weighted sum of career statistics (major titles,
finals, Masters titles, weeks at No.~1, win rate) and peak Elo rating. Elo
rewards beating strong opponents, which title counts miss; every player starts
at 0 and we use $K=32$. A player is eligible five
years after their last match. We choose the score threshold that maximises F1
on players eligible up to \SplitYear{} (threshold \Threshold).

\section{Evaluation}
Table~\ref{tab:metrics} compares the score with the Hall's own automatic-ballot
rule on the \TestPlayers{} contenders who became eligible after \SplitYear{} and by \AsOf{}, of
whom \TestInducted{} were inducted.

\begin{table}[h]
\centering
\MetricsTable
\caption{Agreement with real inductions on held-out years.}
\label{tab:metrics}
\end{table}

\begin{table}[h]
\centering
\small
\TopTable
\caption{The fifteen highest career scores.}
\label{tab:top}
\end{table}

\bibliographystyle{plainnat}
\bibliography{references}

\end{document}
